\usepackage[scr=rsfs,cal=euler]{mathalfa}
\newenvironment{enumeratei}
{\bgroup\def\theenumi{\roman{enumi}}\def\theenumii{\arabic{enumii}}\begin{enumerate}}
{\end{enumerate}\egroup}

\newcommand{\lcr}{[\![}
\newcommand{\rcr}{]\!]}

\newcommand{\Efunction}{$E$\nobreakdash-func\-tion\xspace}
\newcommand{\Efunctions}{$E$\nobreakdash-func\-tions\xspace}

\newtheorem{theo}{Theorem}
\newtheorem{conj}{Conjecture}
\newtheorem{coro}{Corollary}
\newtheorem{prop}{Proposition}
\newtheorem{spec}{Speculation}
\newtheorem{lem}{Lemma}
\newtheorem{defilem}{Definition-Lemma}

\newtheorem{thmx}{Theorem}
\renewcommand{\thethmx}{\Alph{thmx}}

\theoremstyle{definition}
\newtheorem{defi}{Definition}
\newtheorem{Remark}{Remark}
\newtheorem{Remarks}{Remarks}
\newtheorem*{Remark*}{Remark}
\newtheorem*{Remarks*}{Remarks}

\makeatletter
\newcommand*{\house}[1]{%
\mathord{%
\mathpalette\@house{#1}%
}%
}
\newcommand*{\@house}[2]{%
\dimen@=\fontdimen8 %
\ifx#1\scriptscriptstyle\scriptscriptfont
\else\ifx#1\scriptstyle\scriptfont
\else\textfont\fi\fi
3 %
\sbox0{%
$#1%
\vrule width\dimen@\relax
\overline{%
\kern2\dimen@
\begingroup % to keep changes of \dimen@ in #2 local
#2%
\endgroup
\kern2\dimen@
}%
\vrule width\dimen@\relax
\mathsurround=1.5\dimen@ % outside margin
$%
}%
\ht0=\dimexpr\ht0-\dimen@\relax
\dp0=\dimexpr\dp0+2\dimen@\relax
\vbox{%
\kern\dimen@ % reinsert previously removed space
\copy0 %
}%
}
\makeatother

\newcommand{\tra}{{}^t}
\newcommand{\N}{\mathbb{N}}
\newcommand{\Z}{\mathbb{Z}}
\newcommand{\Q}{\mathbb{Q}}
\newcommand{\R}{\mathbb{R}}
\newcommand{\C}{\mathbb{C}}
\newcommand{\K}{\mathbb{K}}
\renewcommand{\L}{\mathbb{L}}

\newcommand{\GG}{{\mathbf G}}
\newcommand{\EE}{{\mathbf E}}

\newcommand{\E}{\EE}
\newcommand{\Qbar}{\overline{\mathbb Q}}
\newcommand{\EK}{\EE_{\K}}
\newcommand{\EQ}{\EE_{\Q}}
\newcommand{\EQbar}{\EE_{\Qbar}}
\newcommand{\Ok}{{\mathcal O}_{\K}}
\newcommand{\OL}{{\mathcal O}_{\L}}
\newcommand{\OQdexi}{{\mathcal O}_{\Q(\xi)}}

\let\eps\varepsilon

\DeclareMathOperator{\Id}{Id}
\DeclareMathOperator{\Gal}{Gal}

\newcommand{\calI}{\mathcal I}
\newcommand{\calE}{\mathcal E}
\newcommand{\calN}{\mathcal N}
\DeclareMathOperator{\Card}{Card}

\newcommand{\isoul}{{\underline{i}}}
\newcommand{\jsoul}{{\underline{j}}}
\newcommand{\nsoul}{\underline{n}}

